\chapter{Wire Formats}
\label[appendix]{app:wire}

This appendix gives the byte layouts that the middleware places on the wire and on disk, read
from the source \citep{scuffedrdma}. Offsets and sizes are in bytes.

\section{Blob Channel Slot}

Each slot begins with the header of \cref{tab:slothdr}, in little-endian
byte order, followed by the key and the chunk (\cref{fig:slot}). The fields
take 40 bytes and the header is padded to 64.

\begin{table}[htbp]
  \centering
  \caption[Header of a blob channel slot]{Header of a blob channel slot.}
  \label{tab:slothdr}
  \small
  \begin{tabular}{L{1.5cm}L{1.3cm}L{3.2cm}L{7.6cm}}
    \toprule
    Offset & Size & Field & Meaning \\
    \midrule
    0 & 4 & magic & \code{0x53424C42}, the letters SBLB; checked on receipt \\
    4 & 4 & blob identifier & counter kept by the sender, the same for every chunk of one blob \\
    8 & 2 & key length & length of the key that follows the header \\
    10 & 2 & flags & bit 0 set if the blob was sent on the hot class \\
    12 & 4 & reserved & zero \\
    16 & 8 & total length & length of the whole blob \\
    24 & 8 & chunk offset & position of this chunk in the blob \\
    32 & 8 & chunk length & length of this chunk \\
    40 & 24 & padding & unused \\
    \bottomrule
  \end{tabular}
\end{table}

The 32-bit immediate value of each write names what arrived. Its top four
bits give the kind, 1 for a data chunk and 2 for a credit, and its low 16
bits give the slot number. A credit is an 8-byte write on the hot endpoint
whose contents are not read; only its immediate value matters.

\section{Packed Quantizer Batch}

A batch of $n$ packed vectors begins with the 20-byte header of
\cref{tab:scqhdr}, in little-endian byte order.

\begin{table}[htbp]
  \centering
  \caption[Header of a packed quantizer batch]{Header of a packed quantizer batch.}
  \label{tab:scqhdr}
  \small
  \begin{tabular}{L{1.5cm}L{1.3cm}L{3.2cm}L{7.6cm}}
    \toprule
    Offset & Size & Field & Meaning \\
    \midrule
    0 & 4 & magic & the letters SCQ1 \\
    4 & 1 & version & 1 \\
    5 & 1 & bits & bits per coordinate, $b$ \\
    6 & 2 & sketch bits & $m$ \\
    8 & 4 & count & number of vectors, $n$ \\
    12 & 4 & dimension & $d$ \\
    16 & 4 & padded dimension & $d$ rounded up to a power of two \\
    \bottomrule
  \end{tabular}
\end{table}

Four arrays follow the header in this order: the $n$ vector lengths as 32-bit
floating point; the $n$ residual lengths as 32-bit floating point; the codes,
$b$ bits each, packed with the lowest bit first; and the sketch signs, one
bit each, packed the same way. The two packed arrays are each rounded up to
a whole byte for the batch, so the size of a batch in bytes is
\mathalt{s batch equals 20 plus 4 n plus 4 n plus the ceiling of n times d times b over 8, plus
  the ceiling of n times m over 8}{%
\begin{equation}
  s_{\text{batch}} = 20 + 4n + 4n + \left\lceil \frac{n\, d\, b}{8} \right\rceil + \left\lceil \frac{n\, m}{8} \right\rceil ,
\end{equation}}
where $20$ is the size of the header, $n$ is the number of vectors, $d$ is
their dimension, with the padded dimension in its place when the two differ,
$b$ is the number of bits per coordinate, and $m$ is the number of sign bits
of the sketch. Per vector and without the shared header this is
\cref{eq:packed}.

\section{Cache Blob for the Connector}

One blob carries the cache blocks of one layer of one request. Its layout is
in \cref{tab:skv}.

\begin{table}[htbp]
  \centering
  \caption[Layout of a cache blob]{Layout of a cache blob.}
  \label{tab:skv}
  \small
  \begin{tabular}{L{1.5cm}L{2.2cm}L{9.9cm}}
    \toprule
    Offset & Size & Content \\
    \midrule
    0 & 4 & the letters SKV1 \\
    4 & 4 & length $h$ of the header, big-endian \\
    8 & $h$ & header, a JSON object in UTF-8 \\
    $8 + h$ & rest & the arrays, one after another \\
    \bottomrule
  \end{tabular}
\end{table}

The header holds a format version, the layer name, the axis along which
blocks are stacked, the block size, the number of tokens, the codec, the
original data type and shape, and one entry per array with its name, data
type, shape, offset and length. Data types travel as names, because the
16-bit format the engine uses has no native type in the array library. With
the raw codec there is one array, the gathered blocks. With the quantizer
codec the arrays are those of the packed batch above.

A control message uses the same first eight bytes with a header length of
zero, followed directly by a small JSON object. The ``done'' message that
ends a request's layers is one of these and travels on the hot class.

\section{Connection Record of the Python Pool}

The Python implementation of the pool on verbs, which produced the loopback measurement of \cref{sec:res-loopback}, exchanges one fixed record per queue pair over TCP before connecting. It is 64 bytes in network byte order
(\cref{tab:qpinfo}). The libfabric pool does not use it; its side channel
exchanges addresses and keys instead.

\begin{table}[htbp]
  \centering
  \caption[Connection record of the Python pool]{Connection record of the Python pool.}
  \label{tab:qpinfo}
  \small
  \begin{tabular}{L{1.5cm}L{1.3cm}L{3.2cm}L{7.6cm}}
    \toprule
    Offset & Size & Field & Meaning \\
    \midrule
    0 & 4 & queue pair number & identifies the queue pair at its owner \\
    4 & 4 & sequence number & first packet sequence number the owner will send \\
    8 & 16 & GID & address of the port on the RDMA network \\
    24 & 4 & IPv4 address & the owner's IP address \\
    28 & 4 & GID index & which entry of the port's address table was used \\
    32 & 4 & MTU & RDMA packet size \\
    36 & 28 & reserved & zero \\
    \bottomrule
  \end{tabular}
\end{table}
